\chapter{Ký hiệu}
\label{sec:notation}

Các loại nơ-ron được ký hiệu bằng bốn chữ cái đầu trong tên tiếng Anh của chất dẫn truyền thần kinh chính: \nt{GLUT} --- glutamate, \nt{GABA} --- axit gamma-aminobutyric, \nt{GLYC} --- glycine, \nt{ACET} --- acetylcholine, \nt{DOPA} --- dopamine, \nt{SERO} --- serotonin, \nt{HIST} --- histamine, \nt{NORA} --- noradrenaline, \nt{ADRE} --- adrenaline, \nt{ASPA} --- aspartate (bảng~\ref{tab:types}).

Chữ cái ít hơn đại lượng, nên ở các chương khác nhau đôi khi cùng một chữ mang nghĩa khác nhau. Trong bảng dưới đây, mỗi nghĩa có một dòng riêng; cột bên phải là công thức nơi nó được đưa ra.

\begin{center}
\footnotesize
\setlength{\tabcolsep}{4pt}
\renewcommand{\arraystretch}{1.12}
\begin{longtable}{>{\raggedright}p{2.0cm}>{\raggedright}p{9.6cm}>{\raggedright\arraybackslash}p{2.4cm}}
\toprule
Ký hiệu & Ý nghĩa & Được đưa ra ở \\
\midrule
\endfirsthead
\toprule
Ký hiệu & Ý nghĩa & Được đưa ra ở \\
\midrule
\endhead
\bottomrule
\endfoot
$a$ & độ dốc của đặc tuyến truyền tuyến tính, $f(u)=au$ & \eqref{eq:feedback} \\
$a$ & mức \nt{ACET}: $0$ --- đọc, $1$ --- ghi & \eqref{eq:acetnet} \\
$a_i$, $a$ & sự mỏi của nhóm: trong bộ tạo nhịp; trong dao động của giấc ngủ sóng chậm & \eqref{eq:halfcentre}, \eqref{eq:updown} \\
$a_S$, $a_P$ & độ nhạy của nhịp tim với chân ga và với phanh & \eqref{eq:gasbrake} \\
$A(r)$, $A_0$ & đáp ứng của khớp thần kinh với một xung ở tần số $r$; của khớp đã nghỉ & \eqref{eq:depression} \\
$A_1$ & đáp ứng của bên nhận với một gói chất dẫn truyền & \eqref{eq:quantal} \\
$A$ & diện tích màng của nơ-ron tính cả đuôi gai & \eqref{eq:dendriteload} \\
$A_f$, $A_s$ & phần hiệu chỉnh được giữ lại đến lần chuyển động sau, ở quá trình nhanh và quá trình chậm & \eqref{eq:tworate} \\
$b_i$ & dòng vào riêng của bộ tạo nhịp & \eqref{eq:seroneuron} \\
$b$ & công việc làm mỏi nhóm nhanh đến đâu, trong bộ tạo nhịp hoặc trong dao động của giấc ngủ & \eqref{eq:halfcentre}, \eqref{eq:updown} \\
$b$ & số bit cho một mẫu số hóa & \eqref{eq:probedata} \\
$B_f$, $B_s$ & mức độ phần hiệu chỉnh học từ sai số, ở quá trình nhanh và quá trình chậm & \eqref{eq:tworate} \\
$c$ & nồng độ chất dẫn truyền trong thể tích & \eqref{eq:seroneuron} \\
$c$ & hệ số tương quan giữa nhiễu của các nơ-ron lân cận & \eqref{eq:correlated} \\
$c$ & hệ số theo đó hiệu số đếm đẩy vợt & \eqref{eq:dishreadout} \\
$c_f$ & đáp ứng của nơ-ron $C$: giá trị lớn nhất của $s_{f,p}$ trên mọi vị trí & \eqref{eq:scstage} \\
$c_m$ & điện dung riêng của màng & \eqref{eq:dendriteload} \\
$C$ & giá của nỗ lực: hành động trong thời gian $\tau_a$ tốn $C/\tau_a$ & \eqref{eq:vigor} \\
$d'$ & độ phân biệt của hai đầu vào: hiệu số chia cho nhiễu & \eqref{eq:gainsnr} \\
$d$, $d_j$ & độ trễ trên đường đi của tín hiệu & \eqref{eq:glutneuron}, \eqref{eq:vstar} \\
$d$ & độ trễ của vòng phản hồi với cơ & \eqref{eq:delaystable} \\
$D$ & hệ số khuếch tán & \eqref{eq:diffusionlength} \\
$D$ & mẫu số trong tần số của mạng \nt{GLUT}--\nt{GABA} & \eqref{eq:isn} \\
$D$ & thời gian chờ phần thưởng bị trì hoãn & \eqref{eq:patience} \\
$e$, $e_{ij}$ & vết ở khớp thần kinh & \eqref{eq:threefactor} \\
$e$, $e_i$ & sai số của đầu ra hoặc của chuyển động, đến qua dây dẫn sai số & \eqref{eq:lms}, \eqref{eq:adaptivefilter}, \eqref{eq:tworate} \\
$E_{\text{xung}}$ & năng lượng của một xung tính cả công việc của khớp thần kinh & \eqref{eq:energy} \\
$E$ & tần số của nhóm \nt{GLUT} & \eqref{eq:ei} \\
$E_E$ & điện thế cân bằng của khớp thần kinh kích thích & \eqref{eq:shunt} \\
$f$, $F$ & đặc tuyến truyền: tần số như một hàm của đầu vào & \eqref{eq:ratenet}, \eqref{eq:gain} \\
$f$, $f_0$ & tần số nhịp tim; nhịp riêng của bộ tạo nhịp & \eqref{eq:gasbrake} \\
$f_s$ & tần số số hóa của một điện cực & \eqref{eq:probedata} \\
$F(t)$, $F_1$ & lực của cơ; lực của một cú giật & \eqref{eq:forcesum} \\
$F_1$, $F_2$ & lực của hai cơ kéo khớp theo hai hướng ngược nhau & \eqref{eq:antagonists} \\
$g$ & độ nhạy của thụ thể với nồng độ & \eqref{eq:seroneuron} \\
$g_L$, $g_E$, $g_I$ & độ dẫn rò, độ dẫn kích thích và độ dẫn ức chế & \eqref{eq:shunt} \\
$g$ & hệ số khuếch đại do \nt{NORA} đặt & \eqref{eq:gain}, \eqref{eq:machine} \\
$g$ & ức chế chung từ \nt{GABA} trong mạng nhớ & \eqref{eq:acetnet} \\
$g$ & hệ số khuếch đại của cổng đau, được đặt từ trên xuống & \eqref{eq:gate} \\
$g$ & hệ số khuếch đại của một tầng trong thác hormone & \eqref{eq:cascade} \\
$g_j$ & bộ lọc $j$ với hằng số thời gian riêng ở đầu vào bộ lọc thích nghi & \eqref{eq:adaptivefilter} \\
$G$ & hệ số khuếch đại của vòng phản hồi & \eqref{eq:feedback} \\
$G$ & tốc độ sửa sai trong vòng có trễ & \eqref{eq:delaystable} \\
$h$, $h_I$ & đầu vào ngoài của nhóm \nt{GLUT} và nhóm \nt{GABA} & \eqref{eq:ei}, \eqref{eq:isn} \\
$h(t)$ & một cú giật đơn của cơ & \eqref{eq:twitch} \\
$H(t)$ & ngưỡng trên của áp lực ngủ: chạm tới đó thì cỗ máy ngủ & \eqref{eq:sleeppressure} \\
$I$, $I_i$ & dòng vào từ bên ngoài tới nơ-ron hoặc nhóm & \eqref{eq:ratenet}, \eqref{eq:wta}, \eqref{eq:updown} \\
$I$ & tần số của nhóm \nt{GABA} & \eqref{eq:ei} \\
$I$ & độ sáng hoặc cường độ của kích thích & \eqref{eq:photonnoise}, \eqref{eq:nakarushton} \\
$I$ & dòng điện đầu vào nạp cho màng & \eqref{eq:dendriteload} \\
$k$ & số xung trong một chùm xung & \eqref{eq:burst} \\
$k$ & phần tăng thêm phải lớn hơn nhiễu bao nhiêu lần & \eqref{eq:photonnoise} \\
$k$ & hệ số phản hồi của thác hormone & \eqref{eq:cascade} \\
$k_s$ & mức độ cơn đau bơm thêm vào sự nhạy cảm hóa & \eqref{eq:sensitization} \\
$K$ & số đặc trưng của kết quả & \eqref{eq:vectortd} \\
$K$ & độ cứng của khớp & \eqref{eq:antagonists} \\
$K$ & hệ số khuếch đại vòng của thác hormone, $K=g^3k$ & \eqref{eq:cascade}, \eqref{eq:cascadestable} \\
$K_\varphi$ & độ mạnh của việc chỉnh bộ tạo nhịp ngày đêm theo tín hiệu ngoài & \eqref{eq:pll} \\
$L$ & độ dài đường dây trễ & \eqref{eq:itd} \\
$L$ & độ dài dây dẫn từ cảm biến đau & \eqref{eq:paintime} \\
$L(t)$ & ngưỡng dưới của áp lực ngủ: chạm tới đó thì cỗ máy thức dậy & \eqref{eq:sleeppressure} \\
$M$ & thước đo mối liên hệ nhân quả giữa tín hiệu báo trước và phần thưởng & \eqref{eq:retro} \\
$M_k$ & kỳ vọng của đặc trưng $k$ & \eqref{eq:vectortd} \\
$M_{ik}$ & độ mạnh ức chế từ nơ-ron \nt{GABA} $k$ & \eqref{eq:machine} \\
$M$ & mô-men quay của khớp & \eqref{eq:antagonists} \\
$M$ & số đường đầu vào của các bộ trộn & \eqref{eq:cover} \\
$n$, $\bar n$ & số xung trong một cửa sổ; giá trị trung bình của nó & \eqref{eq:rateerror} \\
$n$ & số đầu vào của phần tử ngưỡng & \eqref{eq:cover} \\
$n_\uparrow$, $n_\downarrow$ & số đếm xung của các vùng ``lên'' và ``xuống'' trong một cửa sổ & \eqref{eq:dishreadout} \\
$N$ & số nơ-ron & \eqref{eq:correlated}, \eqref{eq:energy} \\
$N$ & số điện cực của đầu dò & \eqref{eq:probedata} \\
$N_s$ & số vị trí nhả giữa hai nơ-ron & \eqref{eq:quantal} \\
$p$ & xác suất nhả chất dẫn truyền ứng với một xung & \eqref{eq:burst} \\
$p$ & tỷ lệ nơ-ron hoạt động trong mạng nhớ & \eqref{eq:acetnet} \\
$p$ & nhiễu loạn mà hệ học cách bù trừ & \eqref{eq:tworate} \\
$P$ & công suất tiêu tốn cho tín hiệu & \eqref{eq:energy}, \eqref{eq:dishpower} \\
$P$ & độ bất định của ước tính đang lưu & \eqref{eq:kalman} \\
$P$ & hoạt động của ``phanh'': dây dẫn \nt{ACET} tới tim & \eqref{eq:gasbrake} \\
$P_{\max}$ & số mẫu ngẫu nhiên tối đa mà phần tử ngưỡng có thể chia thành hai lớp & \eqref{eq:cover} \\
$P_{\text{trượt}}$ & xác suất trượt & \eqref{eq:dishdensity} \\
$q$ & tốc độ thay đổi của thế giới & \eqref{eq:kalman} \\
$q_k$ & tần số của nơ-ron \nt{GABA} $k$ & \eqref{eq:machine} \\
$q_0$ & hoạt động nền của nơ-ron trung gian ức chế & \eqref{eq:gate} \\
$q$ & tần số của nơ-ron trung gian ức chế trong cổng đau & \eqref{eq:gate} \\
$q$ & đơn vị vận động kế tiếp mạnh hơn đơn vị trước bao nhiêu lần & \eqref{eq:sizeprinciple} \\
$r$, $r_i$ & tần số xung của nơ-ron & \eqref{eq:ratenet} \\
$r$, $r_t$ & phần thưởng (dòng phần thưởng) & \eqref{eq:rpe}, \eqref{eq:ctd} \\
$\mathbf r(t)$ & vectơ đáp ứng của hồ chứa & \eqref{eq:reservoir} \\
$R$, $\bar R$ & phần thưởng nhận được; kỳ vọng của nó hoặc giá trị trung bình trên một đơn vị thời gian & \eqref{eq:perturbation}, \eqref{eq:vigor} \\
$R$ & phương sai của nhiễu ở đầu vào bộ lọc & \eqref{eq:kalman} \\
$R$, $r_0$ & phần thưởng bị trì hoãn và phần thưởng tức thì & \eqref{eq:patience} \\
$R_{\max}$ & tốc độ truyền giới hạn, bit/s & \eqref{eq:spikeentropy} \\
$r_{\max}$ & tần số xung lớn nhất & \eqref{eq:gain}, \eqref{eq:nakarushton} \\
$R_{\text{thô}}$, $R_{\text{xung}}$ & luồng dữ liệu của đầu dò: thô và chỉ các xung, bit/s & \eqref{eq:probedata} \\
$s_j$ & dấu của các đầu ra của nơ-ron $j$ & \eqref{eq:dalesign} \\
$s_i$ & dấu của thụ thể ở bên nhận & \eqref{eq:seroneuron} \\
$s$, $\hat s$ & trạng thái của thế giới; ước tính của nó & \eqref{eq:rpe}, \eqref{eq:kalman} \\
$s_{f,p}$ & đáp ứng của nơ-ron $S$ với đặc trưng $f$ ở vị trí $p$ & \eqref{eq:scstage} \\
$S$ & hoạt động của ``chân ga'': dây dẫn \nt{NORA} tới tim & \eqref{eq:gasbrake} \\
$S$ & áp lực ngủ & \eqref{eq:sleeppressure} \\
$t_1$ & thời điểm của xung đầu tiên & \eqref{eq:latency} \\
$t_k$ & thời điểm của xung thứ $k$ & \eqref{eq:forcesum} \\
$t_{f,i}$ & mẫu của đặc trưng $f$ & \eqref{eq:scstage} \\
$T$ & cửa sổ đếm xung; thời gian tích lũy photon & \eqref{eq:rateerror}, \eqref{eq:photonnoise} \\
$T_c$ & thời gian tăng của cú giật cơ & \eqref{eq:twitch} \\
$T_{\text{đợt}}$ & khoảng cách giữa các đợt bùng phát của mẫu nuôi cấy & \eqref{eq:burstinterval} \\
$u$ & tổng đầu vào của nơ-ron & \eqref{eq:gain} \\
$u$, $u(t)$ & lệnh của não: ở đầu vào bộ lọc thích nghi; tới thác hormone & \eqref{eq:adaptivefilter}, \eqref{eq:cascade} \\
$u(t)$ & đầu vào của hồ chứa & \eqref{eq:reservoir} \\
$U$ & mức của đầu vào không đổi & \eqref{eq:latency} \\
$U$ & phần của kho chất dẫn truyền được nhả ra ứng với một xung & \eqref{eq:depression} \\
$v$ & tốc độ tín hiệu trong dây dẫn; tốc độ di chuyển của vệt sáng & \eqref{eq:itd}, \eqref{eq:vstar} \\
$V$ & điện áp trên màng & \eqref{eq:glutneuron}, \eqref{eq:shunt} \\
$V$ & kỳ vọng phần thưởng tương lai & \eqref{eq:rpe}, \eqref{eq:ctd} \\
$w$, $w_{ij}$, $W_{ij}$ & trọng số liên kết & \eqref{eq:glutneuron}, \eqref{eq:ratenet}, \eqref{eq:acetnet}, \eqref{eq:lms}, \eqref{eq:adaptivefilter} \\
$\mathbf w_k$ & trọng số các đầu vào của nơ-ron $C$ thứ $k$ & \eqref{eq:traceinvariance} \\
$w_{q\mu}$, $w_{q\nu}$, $w_{z\mu}$ & trọng số các liên kết trong cổng đau & \eqref{eq:gate} \\
$W_{\text{ra}}$ & trọng số của lớp đọc ra tuyến tính của hồ chứa & \eqref{eq:reservoir} \\
$x$, $y$ & các đầu vào và đầu ra logic & \eqref{eq:dualrail}, \eqref{eq:veto} \\
$x$, $y$ & hoạt động của bên gửi và bên nhận & \eqref{eq:hebb} \\
$x$ & vị trí của bộ phát hiện trên đường dây trễ & \eqref{eq:itd} \\
$x_i$ & đầu vào từ các giác quan & \eqref{eq:acetnet} \\
$x_p$ & đầu vào ở vị trí $p$ & \eqref{eq:scstage} \\
$\mathbf x$ & các đầu ra của nơ-ron $S$, đầu vào của nơ-ron $C$ & \eqref{eq:traceinvariance} \\
$x_j$ & đầu vào $j$ của bộ lọc thích nghi: lệnh sau bộ lọc $g_j$ & \eqref{eq:lms}, \eqref{eq:adaptivefilter} \\
$x_f$, $x_s$ & phần hiệu chỉnh của quá trình nhanh và quá trình chậm & \eqref{eq:tworate} \\
$x_1$, $x_2$, $x_3$ & các mức ở ba tầng của thác hormone; $x_3$ --- hormone cuối cùng & \eqref{eq:cascade} \\
$x^*$ & tỷ lệ kho chất dẫn truyền đã hồi phục mà tại đó một trận thác mới bắt đầu & \eqref{eq:burstinterval} \\
$y$ & đầu vào có nhiễu của bộ lọc & \eqref{eq:kalman} \\
$y$, $\hat y$ & tín hiệu cảm biến; dự đoán của nó bởi bộ lọc thích nghi & \eqref{eq:adaptivefilter} \\
$y_k$, $\bar y_k$ & hoạt động của nơ-ron $C$ thứ $k$; vết của nó & \eqref{eq:traceinvariance} \\
$y(t)$ & đầu ra của lớp đọc ra của hồ chứa & \eqref{eq:reservoir} \\
$z$ & tần số đầu ra của cổng đau & \eqref{eq:gate} \\
\midrule
$\alpha_i^\pm$ & trọng số của sai số dương và sai số âm & \eqref{eq:asymmetric} \\
$\beta$ & độ mạnh của ức chế lẫn nhau & \eqref{eq:wta} \\
$\beta$ & độ mạnh ức chế đầu ra của cổng đau & \eqref{eq:gate} \\
$\gamma$ & hệ số chiết khấu & \eqref{eq:rpe} \\
$\delta(\cdot)$ & hàm delta: bước nhảy tức thời & \eqref{eq:glutneuron} \\
$\delta$, $\delta_t$ & sai số dự đoán phần thưởng & \eqref{eq:rpe}, \eqref{eq:ctd} \\
$\delta t$ & chênh lệch thời điểm âm thanh tới hai tai & \eqref{eq:itd} \\
$\Delta$, $\Delta_{\max}$ & khoảng cách giữa các xung; cửa sổ trùng khớp & \eqref{eq:window} \\
$\Delta t$ & độ chính xác mà bên nhận phân biệt được thời gian & \eqref{eq:spikeentropy} \\
$\Delta F$ & lực tăng thêm từ đơn vị vận động kế tiếp & \eqref{eq:sizeprinciple} \\
$\Delta I$ & phần tăng độ sáng phân biệt được & \eqref{eq:photonnoise} \\
$\Delta p$ & độ dịch của vợt trong một cửa sổ & \eqref{eq:dishreadout} \\
$\Delta u$ & hiệu đầu vào của hai nhóm & \eqref{eq:gainsnr} \\
$\Delta V$ & cần nâng điện áp màng lên bao nhiêu & \eqref{eq:dendriteload} \\
$\Delta x$ & khoảng cách giữa các cảm biến lân cận & \eqref{eq:vstar} \\
$\varepsilon$ & hiệu tương đối của các đầu vào & \eqref{eq:wtagain} \\
$\eta$ & tốc độ học & \eqref{eq:plasticity}, \eqref{eq:lms}, \eqref{eq:traceinvariance} \\
$\eta$ & nhiễu trong mạng sau bộ khuếch đại & \eqref{eq:softmax} \\
$\mu$ & đầu vào xúc giác tới cổng đau & \eqref{eq:gate} \\
$\nu$ & đầu vào đau & \eqref{eq:gate} \\
$\theta$ & ngưỡng kích hoạt & \eqref{eq:window}, \eqref{eq:scstage} \\
$\kappa$ & lượng chất dẫn truyền ứng với một xung & \eqref{eq:seroneuron} \\
$\kappa_i$ & tỷ số giữa trọng số của sai số dương và sai số âm & \eqref{eq:expectile} \\
$\kappa_m$ & liên hệ giữa lực của cơ và độ cứng của nó & \eqref{eq:antagonists} \\
$\ell$ & khoảng cách mà chất dẫn truyền lan ra & \eqref{eq:diffusionlength} \\
$\ell$ & cánh tay đòn mà cơ kéo khớp & \eqref{eq:antagonists} \\
$\xi$ & dao động ngẫu nhiên ở đầu ra của nơ-ron & \eqref{eq:perturbation} \\
$\rho(\boldsymbol\psi)$ & tỷ lệ thời gian ở trong cấu hình $\boldsymbol\psi$ & \eqref{eq:dishdensity} \\
$\sigma$ & độ tản mạn, thang của nhiễu & \eqref{eq:rateerror}, \eqref{eq:softmax} \\
$\sigma$ & độ sáng mà tại đó đáp ứng bằng một nửa giá trị lớn nhất & \eqref{eq:nakarushton} \\
$\sigma_{\text{trước}}$, $\sigma_{\text{sau}}$ & nhiễu trước bộ khuếch đại và sau nó & \eqref{eq:gainsnr} \\
$\tau$ & hằng số thời gian của màng & \eqref{eq:glutneuron} \\
$\tau$ & hằng số thời gian của một tầng trong thác hormone & \eqref{eq:cascade} \\
$\tau_{\text{hd}}$ & hằng số thời gian của nhóm tự kích thích chính nó & \eqref{eq:perfectint} \\
$\tau_c$ & thời gian sống của chất dẫn truyền trong thể tích & \eqref{eq:seroneuron} \\
$\tau_E$, $\tau_I$ & hằng số thời gian của nhóm \nt{GLUT} và nhóm \nt{GABA} & \eqref{eq:ei} \\
$\tau_e$ & thời gian sống của vết & \eqref{eq:threefactor} \\
$\tau_{\text{tr}}$ & thời gian sống của vết hoạt động của nơ-ron $C$ & \eqref{eq:traceinvariance} \\
$\tau_s$ & thời gian phục hồi độ nhạy sau tổn thương & \eqref{eq:sensitization} \\
$\tau_r$, $\tau_d$ & thời gian tích lũy áp lực ngủ khi thức và thời gian xả nó khi ngủ & \eqref{eq:sleeppressure} \\
$\tau_{\text{ph}}$ & thời gian phục hồi kho chất dẫn truyền & \eqref{eq:depression} \\
$\tau_\gamma$ & tầm nhìn lập kế hoạch & \eqref{eq:ctd} \\
$\tau_v$ & tốc độ chính xác hóa kỳ vọng & \eqref{eq:ramp} \\
$\tau_a$ & thời gian mỏi trong bộ tạo nhịp hoặc trong dao động của giấc ngủ & \eqref{eq:halfcentre}, \eqref{eq:updown} \\
$\tau_a^*$ & thời gian thực hiện hành động tốt nhất & \eqref{eq:vigor} \\
$\Phi$ & vế phải của quy tắc học & \eqref{eq:plasticity} \\
$\Phi$ & dòng photon & \eqref{eq:photonnoise} \\
$\varphi_k$ & đặc trưng $k$ của kết quả nhận được & \eqref{eq:vectortd} \\
$\varphi$ & độ lệch pha giữa bộ tạo nhịp ngày đêm và tín hiệu ngoài & \eqref{eq:pll} \\
$\boldsymbol\psi$ & cấu hình của mạng trong mô hình DishBrain & \eqref{eq:dishdensity} \\
$\omega$, $\omega_0$ & tần số của tín hiệu ngoài (ánh sáng); tần số riêng của bộ tạo nhịp ngày đêm & \eqref{eq:pll} \\
\end{longtable}
\end{center}
